\documentclass[11pt]{article}

% Change "review" to "final" to generate the final (sometimes called camera-ready) version.
% Change to "preprint" to generate a non-anonymous version with page numbers.
\usepackage[]{acl}
\usepackage{times}
\usepackage{latexsym}

% For proper rendering and hyphenation of words containing Latin characters (including in bib files)
\usepackage[T1]{fontenc}
% For Vietnamese characters
% \usepackage[T5]{fontenc}
% See https://www.latex-project.org/help/documentation/encguide.pdf for other character sets

% This assumes your files are encoded as UTF8
\usepackage[utf8]{inputenc}
\usepackage{microtype}
\usepackage{inconsolata}
\usepackage{graphicx}

% If the title and author information does not fit in the area allocated, uncomment the following
%
%\setlength\titlebox{<dim>}
%
% and set <dim> to something 5cm or larger.

\title{Group X Progress Report:\\My Group's Project Name}


\author{Nehan Munir, Momina , Third Author \\
  \texttt{\{munirn3,macid2,macid3\}@mcmaster.ca} }

%\author{
%  \textbf{First Author\textsuperscript{1}},
%  \textbf{Second Author\textsuperscript{1,2}},
%  \textbf{Third T. Author\textsuperscript{1}},
%  \textbf{Fourth Author\textsuperscript{1}},
%\\
%  \textbf{Fifth Author\textsuperscript{1,2}},
%  \textbf{Sixth Author\textsuperscript{1}},
%  \textbf{Seventh Author\textsuperscript{1}},
%  \textbf{Eighth Author \textsuperscript{1,2,3,4}},
%\\
%  \textbf{Ninth Author\textsuperscript{1}},
%  \textbf{Tenth Author\textsuperscript{1}},
%  \textbf{Eleventh E. Author\textsuperscript{1,2,3,4,5}},
%  \textbf{Twelfth Author\textsuperscript{1}},
%\\
%  \textbf{Thirteenth Author\textsuperscript{3}},
%  \textbf{Fourteenth F. Author\textsuperscript{2,4}},
%  \textbf{Fifteenth Author\textsuperscript{1}},
%  \textbf{Sixteenth Author\textsuperscript{1}},
%\\
%  \textbf{Seventeenth S. Author\textsuperscript{4,5}},
%  \textbf{Eighteenth Author\textsuperscript{3,4}},
%  \textbf{Nineteenth N. Author\textsuperscript{2,5}},
%  \textbf{Twentieth Author\textsuperscript{1}}
%\\
%\\
%  \textsuperscript{1}Affiliation 1,
%  \textsuperscript{2}Affiliation 2,
%  \textsuperscript{3}Affiliation 3,
%  \textsuperscript{4}Affiliation 4,
%  \textsuperscript{5}Affiliation 5
%\\
%  \small{
%    \textbf{Correspondence:} \href{mailto:email@domain}{email@domain}
%  }
%}

\begin{document}
\maketitle
% \begin{abstract}
% \end{abstract}

\section{Introduction}

Our project investigates the following question: 

\textit{Which acoustic characteristics are most useful for distinguishing Parkinson's disease from healthy controls, and which machine learning approach best captures these differences?}

Parkinson’s disease (PD) is a progressive neurodegenerative disorder that damages dopamine producing neurons in the brain and affects motor functions as well as speech production.


Parkinson’s is a hard disease to diagnose without either a specialized neurologist or expensive tests, and this results in it being misdiagnosed 30\% of the time (NCBI/NIH, 2020). The main symptom that the team is looking for is Hypokinetic Dysarthria, which affects 89\% of all people with PD (Byatt, 2025). Hypokinetic Dysarthria (HD) is described as “monotonous speech pattern, reduced and monotonous loudness, decreased stress, a breathy or hoarse voice quality, an increase in speech rate, rapid repetition of phonemes, and impreciseness in consonant production”, which can be tested using the dataset that we have procured by testing for the following variables: “phonation, articulation, respiration, resonance, and prosody” (Atalar et al., 2023).


In order for a doctor to diagnose someone as having PD or to ask for further testing, they have to prove that the patient has Bradykinesia alongside tremors/rigidity (Postuma et al., 2015). Bradykinesia is “slowed movement” and it is extremely difficult to diagnose because it is hard to differentiate it from the side effects of aging (Cleveland Clinic, 2023). However, Bradykinesia causes HD, so having a reliable and accurate method to diagnose it will firstly help caretakers/suspect cases realize that they might have PD without paying for a doctor’s visit, and secondly help doctors in asking for a referral for a neurologist/tests. 

\section{Related Work}

Here, talk about the related work you encountered for your approach. Cite at least 5 references. Refer to item 2. No one has done exactly your task? Write about the most similar thing you can find. This should be around 0.25-0.5 pages.

\section{Dataset}

You should write about your dataset here, following the guidelines regarding item 1. This section may be 0.5-1 pages. Depending on your specific dataset, you may want to include subsections for the preprocessing, annotation, etc.

The dataset the team decided to select is Mobile Device Voice Recordings (Trivedi, 2020). These are audio recordings of
37 participants collected by Dhaval Trivedi, Hagen Jaeger and Michael Stadtschnitzer at the King's College Hospital in 
London, UK. The dataset is going to be referred as KCL going forward. Each participant was taken to a 10 m² room where
they were recorded on a Motorola Moto G4 using a self-developed app similar to i-PROGNOSIS. In between May 17, 2019 and January 24, 2020,
the KCL team recorded the participants reading “The North Wind and the Sun” and “Tech. Engin. Computer applications in geography snippet”
and adding them to a "Read Text" folder. They also recorded a simple phone conversation with each participant (except for participant 18),
which were then added to a "Spontaneous Dialogue". They further seperated each folder by their Parkinson's Disease
diagnosis, labelling the sub-folders "PD" (Parkinson's Disease) or "HC" (Human Control). They labelled all .wav files 
with the participant ID, Hoehn and Yahr Scale score, UPDRS II-5 and UPDRS III-18 score, resulting in the following for 
patient 32 (with Parkinson's): ID32\_pd\_3\_1\_1. 

Before trying to find acoustic features, the team looked into seeing whether the numerical features in the labelled .wav files
contain any information. Unfortunately, they are usually used to categorize the progression of PD, which means any number
outside of 0 means that the user already has the disease. 

As stated previously, there are only 37 patients, meaning that there will be a maximum of 74 recordings if every participant
did both aspects of the data gathering. This falls below the recommended dataset size of 1000 data points, which resulted
in the team devising a plan to seperate the approximately 2 minute recordings into 10 seconds segments. This will give 
the team closer to 900 data points. There might be problems with overfitting as the final machine learning model will see
the same voice multiple times, but this can be discouraged based on how the team splits the training data. In order to
preprocess the data, the team creates a better folder directory structure splitted firstly by their diagnosis, and then 
each participant folder (labelled by participant ID) will have subfolders for their Spontaneous Dialogue (SD) and 
Read Text (RT) sections (as they contain the .wav splits). 

The wav files are then created into AudioSegment objects using Pydub, which allows splitting based on milliseconds. Each
wav file has been splitted into a list of increments of 10 seconds, with the last element of the list being the remaining
segment. The last element is discarded if it isn't bigger than 8 seconds (an arbitary number brainstormed as the minimum
time needed to distill information). Each participant has been assigned a dictionary, with the numerical features from the 
labelled wav files (Hoehn and Yahr Scale score, UPDRS II-5 and UPDRS III-18 score) alongside a boolean describing their
PD diagnosis. The splitted wav files are then added to the correspondent folder, following HC/PD -> ID -> RT/SD folder. The
new wav files are labelled starting from 0 (as num.wav). The links to the folders for the RT and SD splits are then
added to the correct participant dictionary. 



\section{Features}

Describe any features you used for your model, or how your data was input to your model. Are you doing feature engineering or feature selection? Are you learning embeddings? Is it all part of one neural network? Refer to item 2. This may range from 0.25 pages to 0.5 pages.

\section{Implementation}

Describe your model and implementation here. Refer to item 4. This may take around a page.

\section{Results and Evaluation}

How are you evaluating your model? What results do you have so far? What are your baselines? Refer to item 5. This may take around 0.5 pages.

\section{Feedback and Plans}

Write about your plans for the remainder of the project. This should include a discussion of the feedback you received from your TA, and how you plan to improve your approach. Reflect on your implementation and areas for improvement. Refer to item 6. This may be around 0.5 pages.

\section{Template Notes}

You can remove this section or comment it out, as it only contains instructions for how to use this template. You may use subsections in your document as you find appropriate.

\subsection{Tables and figures}

See Table~\ref{citation-guide} for an example of a table and its caption.
See Figure~\ref{fig:experiments} for an example of a figure and its caption.


\begin{figure}[t]
  \includegraphics[width=\columnwidth]{example-image-golden}
  \caption{A figure with a caption that runs for more than one line.
    Example image is usually available through the \texttt{mwe} package
    without even mentioning it in the preamble.}
  \label{fig:experiments}
\end{figure}

\begin{figure*}[t]
  \includegraphics[width=0.48\linewidth]{example-image-a} \hfill
  \includegraphics[width=0.48\linewidth]{example-image-b}
  \caption {A minimal working example to demonstrate how to place
    two images side-by-side.}
\end{figure*}


\subsection{Citations}

\begin{table*}
  \centering
  \begin{tabular}{lll}
    \hline
    \textbf{Output}           & \textbf{natbib command} & \textbf{ACL only command} \\
    \hline
    \citep{Gusfield:97}       & \verb|\citep|           &                           \\
    \citealp{Gusfield:97}     & \verb|\citealp|         &                           \\
    \citet{Gusfield:97}       & \verb|\citet|           &                           \\
    \citeyearpar{Gusfield:97} & \verb|\citeyearpar|     &                           \\
    \citeposs{Gusfield:97}    &                         & \verb|\citeposs|          \\
    \hline
  \end{tabular}
  \caption{\label{citation-guide}
    Citation commands supported by the style file.
  }
\end{table*}

Table~\ref{citation-guide} shows the syntax supported by the style files.
We encourage you to use the natbib styles.
You can use the command \verb|\citet| (cite in text) to get ``author (year)'' citations, like this citation to a paper by \citet{Gusfield:97}.
You can use the command \verb|\citep| (cite in parentheses) to get ``(author, year)'' citations \citep{Gusfield:97}.
You can use the command \verb|\citealp| (alternative cite without parentheses) to get ``author, year'' citations, which is useful for using citations within parentheses (e.g. \citealp{Gusfield:97}).

\subsection{References}

\nocite{Ando2005,andrew2007scalable,rasooli-tetrault-2015}

Many websites where you can find academic papers also allow you to export a bib file for citation or bib formatted entry. Copy this into the \texttt{custom.bib} and you will be able to cite the paper in the \LaTeX{}. You can remove the example entries.

\subsection{Equations}

An example equation is shown below:
\begin{equation}
  \label{eq:example}
  A = \pi r^2
\end{equation}

Labels for equation numbers, sections, subsections, figures and tables
are all defined with the \verb|\label{label}| command and cross references
to them are made with the \verb|\ref{label}| command.
This an example cross-reference to Equation~\ref{eq:example}. You can also write equations inline, like this: $A=\pi r^2$.


% \section*{Limitations}

\section*{Team Contributions}

Write in this section a few sentences describing the contributions of each team member. What did each member work on? Refer to item 7.

% Bibliography entries for the entire Anthology, followed by custom entries
%\bibliography{custom,anthology-overleaf-1,anthology-overleaf-2}

% Custom bibliography entries only
\bibliography{custom}

% \appendix

% \section{Example Appendix}
% \label{sec:appendix}

% This is an appendix.

\end{document}
